\def\tituloconcepto{Prerrequisitos y notación}
% Si el capitulo se compila solo (sin main.tex), esta cabecera abre el documento.
% Antes de % Si el capitulo se compila solo (sin main.tex), esta cabecera abre el documento.
% Antes de % Si el capitulo se compila solo (sin main.tex), esta cabecera abre el documento.
% Antes de \input{cabecera} cada capitulo define \tituloconcepto.
\ifdefined\enmain\else
  \documentclass[11pt,a4paper]{article}
  \input{preambulo}
  \begin{document}
  \title{\tituloconcepto}
  \author{\autor}
  \date{\today}
  \maketitle
\fi
 cada capitulo define \tituloconcepto.
\ifdefined\enmain\else
  \documentclass[11pt,a4paper]{article}
  % Preambulo compartido por todos los documentos. Compilar con XeLaTeX/LuaLaTeX
% (o Tectonic). Si se usa Overleaf, elegir el compilador XeLaTeX.
\usepackage[spanish,es-noshorthands]{babel}
\usepackage{amsmath,amssymb}
\usepackage{booktabs,array}
\usepackage[a4paper,margin=2.5cm]{geometry}
\usepackage{xcolor}
\usepackage{listings}
\usepackage[most]{tcolorbox}
\usepackage{quantikz}
\usepackage[hidelinks]{hyperref}

% ---- Autor (editar a gusto) -------------------------------------------------
\newcommand{\autor}{Genesis Silesky}

% ---- Notacion de Dirac y atajos ---------------------------------------------
% quantikz ya define \ket y \bra; DeclareRobustCommand los redefine sin error.
\DeclareRobustCommand{\ket}[1]{\left|#1\right\rangle}
\DeclareRobustCommand{\bra}[1]{\left\langle#1\right|}
\DeclareRobustCommand{\braket}[2]{\left\langle#1\middle|#2\right\rangle}
\newcommand{\C}{\mathbb{C}}
\newcommand{\R}{\mathbb{R}}
\DeclareMathOperator{\Tr}{Tr}

% ---- Cajas de apoyo --------------------------------------------------------
% "concepto": definiciones previas que hacen falta para entender el capitulo.
\newtcolorbox{concepto}[1]{breakable,colback=blue!4,colframe=blue!55!black,
  fonttitle=\bfseries,title={Concepto previo: #1}}
% "nota": aclaraciones y advertencias.
\newtcolorbox{nota}[1]{breakable,colback=orange!6,colframe=orange!70!black,
  fonttitle=\bfseries,title={#1}}

% ---- Estilos de listados ----------------------------------------------------
\lstdefinestyle{py}{language=Python,basicstyle=\ttfamily\footnotesize,
  breaklines=true,postbreak=\mbox{\textcolor{gray}{$\hookrightarrow$}\space},
  numbers=left,numberstyle=\tiny\color{gray},frame=single,
  commentstyle=\color{green!40!black}\itshape,
  keywordstyle=\color{blue!70!black}\bfseries,
  stringstyle=\color{red!60!black},showstringspaces=false,tabsize=4,
  columns=fullflexible,keepspaces=true}
\lstdefinestyle{salida}{basicstyle=\ttfamily\footnotesize,breaklines=true,
  frame=leftline,backgroundcolor=\color{gray!8},columns=fullflexible,
  keepspaces=true}

  \begin{document}
  \title{\tituloconcepto}
  \author{\autor}
  \date{\today}
  \maketitle
\fi
 cada capitulo define \tituloconcepto.
\ifdefined\enmain\else
  \documentclass[11pt,a4paper]{article}
  % Preambulo compartido por todos los documentos. Compilar con XeLaTeX/LuaLaTeX
% (o Tectonic). Si se usa Overleaf, elegir el compilador XeLaTeX.
\usepackage[spanish,es-noshorthands]{babel}
\usepackage{amsmath,amssymb}
\usepackage{booktabs,array}
\usepackage[a4paper,margin=2.5cm]{geometry}
\usepackage{xcolor}
\usepackage{listings}
\usepackage[most]{tcolorbox}
\usepackage{quantikz}
\usepackage[hidelinks]{hyperref}

% ---- Autor (editar a gusto) -------------------------------------------------
\newcommand{\autor}{Genesis Silesky}

% ---- Notacion de Dirac y atajos ---------------------------------------------
% quantikz ya define \ket y \bra; DeclareRobustCommand los redefine sin error.
\DeclareRobustCommand{\ket}[1]{\left|#1\right\rangle}
\DeclareRobustCommand{\bra}[1]{\left\langle#1\right|}
\DeclareRobustCommand{\braket}[2]{\left\langle#1\middle|#2\right\rangle}
\newcommand{\C}{\mathbb{C}}
\newcommand{\R}{\mathbb{R}}
\DeclareMathOperator{\Tr}{Tr}

% ---- Cajas de apoyo --------------------------------------------------------
% "concepto": definiciones previas que hacen falta para entender el capitulo.
\newtcolorbox{concepto}[1]{breakable,colback=blue!4,colframe=blue!55!black,
  fonttitle=\bfseries,title={Concepto previo: #1}}
% "nota": aclaraciones y advertencias.
\newtcolorbox{nota}[1]{breakable,colback=orange!6,colframe=orange!70!black,
  fonttitle=\bfseries,title={#1}}

% ---- Estilos de listados ----------------------------------------------------
\lstdefinestyle{py}{language=Python,basicstyle=\ttfamily\footnotesize,
  breaklines=true,postbreak=\mbox{\textcolor{gray}{$\hookrightarrow$}\space},
  numbers=left,numberstyle=\tiny\color{gray},frame=single,
  commentstyle=\color{green!40!black}\itshape,
  keywordstyle=\color{blue!70!black}\bfseries,
  stringstyle=\color{red!60!black},showstringspaces=false,tabsize=4,
  columns=fullflexible,keepspaces=true}
\lstdefinestyle{salida}{basicstyle=\ttfamily\footnotesize,breaklines=true,
  frame=leftline,backgroundcolor=\color{gray!8},columns=fullflexible,
  keepspaces=true}

  \begin{document}
  \title{\tituloconcepto}
  \author{\autor}
  \date{\today}
  \maketitle
\fi


\section{Prerrequisitos y notación}
\label{sec:prerreq}

Este capítulo reúne las herramientas matemáticas y de software que usan todos
los demás. Cada capítulo de concepto indica, en sus cajas azules, qué
secciones de aquí necesita; si algo no se entiende, conviene volver a la
sección correspondiente antes de continuar.

\subsection{Números complejos}
Un número complejo es $z=a+bi$ con $a,b\in\R$ e $i^2=-1$.
\begin{itemize}
  \item \textbf{Conjugado:} $\bar z=a-bi$.
  \item \textbf{Módulo:} $|z|=\sqrt{z\bar z}=\sqrt{a^2+b^2}$. Es el que importa
    para las probabilidades.
  \item \textbf{Forma polar y fórmula de Euler:} $z=|z|e^{i\varphi}$, con
    $e^{i\varphi}=\cos\varphi+i\sin\varphi$. El ángulo $\varphi$ es la
    \emph{fase}.
\end{itemize}
En Python un complejo se escribe \texttt{1+1j} (la unidad imaginaria es
\texttt{1j}) y \texttt{np.abs(z)} devuelve su módulo.

\subsection{Vectores, norma y producto interno}
Los estados de $n$ qubits viven en $\C^{2^n}$. Para vectores columna
$u,v\in\C^d$ el producto interno es $\braket{u}{v}=\sum_k \bar u_k v_k$ y la
norma es $\|v\|=\sqrt{\braket{v}{v}}$. Un estado cuántico válido siempre tiene
$\|v\|=1$ (``normalizado'').

\subsection{Notación de Dirac}
\begin{center}
\begin{tabular}{@{}lll@{}}
\toprule
Símbolo & Significado & En Python/numpy\\
\midrule
$\ket{\psi}$ & vector columna (``ket'') & \texttt{np.array([a, b])}\\
$\bra{\psi}$ & su transpuesta conjugada (``bra'') & \texttt{v.conj()}\\
$\braket{\varphi}{\psi}$ & producto interno & \texttt{np.vdot(phi, psi)}\\
$\ket{0}=\binom10,\ \ket{1}=\binom01$ & base computacional &
  \texttt{Statevector.from\_label("0")}\\
\bottomrule
\end{tabular}
\end{center}

\subsection{Matrices unitarias}
La adjunta de $U$ es $U^\dagger=\bar U^{\,T}$ (transpuesta conjugada). $U$ es
\emph{unitaria} si $U^\dagger U=UU^\dagger=I$. Las matrices unitarias
conservan la norma: $\|U v\|=\|v\|$, y son invertibles con $U^{-1}=U^\dagger$.
Toda puerta cuántica es una matriz unitaria.

\subsection{Producto tensorial (Kronecker)}
Para combinar sistemas se usa $A\otimes B$. En vectores,
$\ket{a}\otimes\ket{b}$ se escribe $\ket{ab}$; por ejemplo
$\ket{0}\otimes\ket{1}=\ket{01}=(0,1,0,0)^T$. Con $n$ qubits la base tiene
$2^n$ vectores $\ket{x}$, $x\in\{0,1\}^n$. Un estado que \emph{no} puede
escribirse como $\ket{a}\otimes\ket{b}$ está entrelazado (capítulo~05b).

\subsection{Probabilidad básica}
\begin{itemize}
  \item \textbf{Regla de Born:} si el estado es $\sum_x c_x\ket{x}$, medir da
    $x$ con probabilidad $|c_x|^2$ (y esas probabilidades suman 1 por la
    normalización).
  \item \textbf{Frecuencia y ley de los grandes números:} al repetir $N$ veces
    un experimento con probabilidad $p$, la frecuencia observada tiene error
    típico $\sqrt{p(1-p)/N}$ y tiende a $p$.
\end{itemize}

\subsection{Conceptos de Qiskit}
\begin{description}
  \item[Circuito (\texttt{QuantumCircuit})] Lista ordenada de puertas y
    mediciones sobre qubits y bits clásicos.
  \item[\texttt{Statevector}] Calcula el vector de estado \emph{exacto} de un
    circuito sin medir; solo es posible en simulación.
  \item[Simulador (\texttt{AerSimulator})] Ejecuta circuitos en el
    computador local, con o sin ruido.
  \item[Shot] Una ejecución completa del circuito con su medición. Se repite
    muchas veces para estimar probabilidades.
  \item[Conteos (\texttt{get\_counts})] Diccionario cadena-de-bits $\to$ número
    de veces que salió.
\end{description}

\begin{nota}{Orden de los bits en Qiskit (little-endian)}
Qiskit escribe las cadenas con el qubit $0$ a la \textbf{derecha}: la cadena
\texttt{'10'} significa $q_1=1$, $q_0=0$. En estos ejemplos se señala cada
vez que importa.
\end{nota}

\subsection{Puertas usadas en estos ejemplos}
\begin{center}
\begin{tabular}{@{}llll@{}}
\toprule
Puerta & Matriz & Efecto & Método Qiskit\\
\midrule
$H$ & $\frac1{\sqrt2}\begin{psmallmatrix}1&1\\1&-1\end{psmallmatrix}$ &
  superposición & \texttt{qc.h(q)}\\
$X$ & $\begin{psmallmatrix}0&1\\1&0\end{psmallmatrix}$ &
  NOT: $\ket0\leftrightarrow\ket1$ & \texttt{qc.x(q)}\\
$Z$ & $\begin{psmallmatrix}1&0\\0&-1\end{psmallmatrix}$ &
  cambia el signo de $\ket1$ & \texttt{qc.z(q)}\\
$R_y(\theta)$ &
  $\begin{psmallmatrix}\cos\frac\theta2&-\sin\frac\theta2\\
  \sin\frac\theta2&\cos\frac\theta2\end{psmallmatrix}$ &
  rotación real & \texttt{qc.ry(theta, q)}\\
CNOT & (4$\times$4) & invierte el objetivo si el control es $\ket1$ &
  \texttt{qc.cx(c, t)}\\
\bottomrule
\end{tabular}
\end{center}

\subsection{Cómo ejecutar el código}
Cada ejemplo existe como \texttt{.py} (terminal) y como \texttt{.ipynb}
(Jupyter/Colab). En el \texttt{.py}, las líneas \texttt{\# \%\%} separan
``celdas'': en un notebook cada una es una celda; en la terminal no tienen
efecto. Instalación (versiones con las que se probó todo):
\begin{lstlisting}[style=salida]
pip install qiskit==2.5.0 qiskit-aer==0.17.2 matplotlib pylatexenc
python 01_qubit_vector.py
\end{lstlisting}
La API de Qiskit ha cambiado bastante entre versiones; si se usa otra versión
algún ejemplo podría fallar.

\ifdefined\enmain\else
  \end{document}
\fi

